Nach dem Energiesatz ist die kinetische Energie gleich der Abnahme der potentiellen Energie. Für den Aufprall und für die gesuchte Höhe $h_1$ gilt
\begin{align}
 \frac{1}{2}m\,v_\mathrm{A}^2 &= m\,g\,h \\
 \frac{1}{2}m\,v_1^2 &= m\,g\left(h-h_1\right)
\end{align}
Mit $v_1 = v_\mathrm{A}/2$ ist $v_1^2 = v_\mathrm{A}^2/4$. Teilt man die zweite Gleichung durch die erste, folgt
\begin{align}
 \frac{1}{4} &= \frac{h-h_1}{h} \quad\Rightarrow\quad h_1 = \result{\frac{3}{4}h}
\end{align}
Der Stein hat die halbe Aufprallgeschwindigkeit, wenn er erst einen Viertel der Fallhöhe zurückgelegt hat.
